
In this section, we present a nondeterministic polynomial time algorithm for the complement of RIT. As discussed in \Cref{sec:overview}, the idea is to work in a finite field obtained by quotienting the ring of integers of the splitting field of the input radicals by a suitable prime ideal.

Below, we fix an instance of  RIT  given by an algebraic circuit $C$, and input radicals~$\mysqrt{d_1}{a_1},\ldots,\mysqrt{d_k}{a_k}$ with respective minimal polynomials $x^{d_i}-a_i$, where the radicands~$a_i$ are pairwise coprime;
this assumption is without loss of generality as discussed in \cref{appendix:reduction}.
We denote by $s$ the size of our fixed RIT instance; that is, the size of the circuit is bounded by $s$,  and the magnitude of the $a_i$ and $d_i$ is at most $2^s$. Note that  $k\leq s$, by the definition of size of an algebraic circuit. 
 We further  denote by $K$ the splitting field  of $\prod_{i=1}^k (x^{d_i}-a_i)$, which can be generated by adjoining to~$\QQ$ the radicals $\mysqrt{d_i}{a_i}$ and a primitive $d$-th root of unity~$\zeta_d$, with $d = \lcm(d_1,\ldots,d_k)$. We denote by $\OO_K$ the ring of integers of $K$.


In our construction, we evaluate the polynomial given by an algebraic circuit in a finite field $\FF_p$ for some rational prime $p$ that splits completely in $\OO_K$, that is, such that $p\OO_K = \mathfrak{p}_1\cdots \mathfrak{p}_n$ where $\mf{p}_i$ are distinct prime ideals of $\OO_K$, and $n$ is the degree of the number field $K$.

In general, given a number field $L$ and a rational prime $q$ with prime ideal $\mathfrak{q}$ dividing $q\OO_L$, we say that $\mathfrak{q}$ lies above $q$ in $\OO_L$. The residue field $\OO_L/\mathfrak{q}$ is isomorphic to an extension of the finite field~$\FF_q$, 
and we have that $\mathfrak{q} \cap \ZZ=q\ZZ$.
 However, in our special case of a completely split prime $p$ in $\OO_K$, all residue fields $\OO_K/\mathfrak{p}_i$ are isomorphic to $\FF_p$. This is crucial, as it ensures that the values in our finite field computation really will all be in $\FF_p$ and not in one of its finite extensions (see the discussion in \Cref{sec:discussion}).

The following proposition asserts that a prime $p$ completely splits in $K$ if the minimal polynomials $x^{d_i}-a_i$ of our input radicals and the $d$-th cyclotomic polynomial (the minimal polynomial of $\zeta_d$) split into distinct linear factors in $\FF_p$.

\begin{proposition}
	\label{prop:split_primes_qp}
	Given a monic polynomial $g \in\ZZ[x]$ and its splitting field~$L$, 
	a prime $q\in\ZZ$ splits completely in $L$ if $g$ splits into distinct linear factors in $\FF_q$. 
\end{proposition}

\begin{proof}[Proof sketch]
Denote by $L_{\mathfrak{q}}$ the finite field extension of $\QQ_q$ obtained by adjoining the roots of $g$ to $\QQ_q$. Let $\mathfrak{q}$ be a prime ideal lying above $q$ in $L$.
	Recall that the decomposition group of a prime ideal $\mathfrak{q} \subset L$ is defined as the set of all automorphisms of $\Gal(L/\QQ)$ that fix $\mathfrak{q}$, i.e. $D_{\mf{q}} = \{ \sigma \in \Gal(L/\QQ) \mid \sigma(\mathfrak{q}) = \mathfrak{q}\}$. Furthermore, since the field $L$ is Galois, the following isomorphism holds:
\begin{equation}
\label{eq:decomposition_group_iso}
	D_{\mathfrak{q}} \cong \Gal(L_{\mathfrak{q}}/\QQ_q)
\end{equation}

Now given that $g$ splits into distinct linear factors in $\FF_q$, by virtue of $L$ being Galois, it follows that all roots of $g$ in $\FF_q$ are non-zero. Hence we can use Hensel's lemma to construct the solutions of $g$ in $\QQ_q$, that is, $g$ splits completely in $\QQ_q$. This, in turn, implies that $L_{\mathfrak{q}} = \QQ_q$, that is, $\Gal(L_{\mathfrak{q}}/\QQ_q)$ is trivial and \eqref{eq:decomposition_group_iso} asserts that the same holds for $D_{\mathfrak{q}}$. This entails that the only automorphism that fixes the prime ideal $\mathfrak{q}$ is the identity, whereas for all non-trivial $\sigma\in\Gal(L/\QQ)$, which permute the conjugates of our input radicals, they will map elements of $\mathfrak{q}$ to elements of $\OO_L/\mathfrak{q}\cong \FF_q$, i.e., $\FF_q$ will indeed contain all roots of $g$.
\end{proof}

\begin{example}
Consider an RIT instance asking whether the polynomial $f(x) = x^2 - 10$ vanishes at the radical input $\sqrt{5}$. The computation occurs in the field $L=\QQ(\sqrt{5})$, with ring of integers $\OO_L = \ZZ[\frac{1+\sqrt{5}}{2}]$.
	We can observe that $11$ is a completely split prime and note that the principal ideal of $\OO_L$ generated by~$11$ factors as $11\OO_L = (4+\sqrt{5}) (4-\sqrt{5})$.
Since  $11$ totally splits in~$\OO_L$, we have $\OO_L/\mathfrak{q} \cong \FF_{11}$ for the prime ideal $\mathfrak{q}=(4+\sqrt{5})$ lying above~$11$.
The polynomial $x^2-5$  completely splits to  $(x-4)(x+4) $ in $\FF_{11}$, and    consequently we have that $\QQ_{11}(\sqrt{5}) = \QQ_{11}$. 
	The rational prime~$5$, however, is an example of the primes that we want to avoid as $5\OO_L = (\sqrt{5})^2$.
	 In particular, the polynomial $x^{2}-5$ is irreducible in~$\FF_{5}$, implying that $L_{\mathfrak{q}} = \QQ_q(\sqrt{5})$ and $[L_{\mathfrak{q}} : \QQ_q] = 2$. If we evaluate $f$ on the input $\sqrt{5}$ in $\FF_{11}$, we get a true negative, as the computed value will be $6\in\FF_{11}$, whereas evaluating the polynomial in $\FF_5$ would give us a false positive.
\end{example}


\subsection{Proof of correctness}
\label{sec:correctness}

Given a polynomial~$g\in \ZZ[x]$ that is irreducible over~$\QQ$, the Galois group $\Gal(L/\QQ)$ of the splitting field~$L$ of $g$ acts transitively on the roots of $g$~\cite[Proposition~22.3]{StewartBook} . We show that a stronger notion of transitivity holds for our real radicals~$\mysqrt{d_i}{a_i}$:  

\begin{lemma}\label{lem:jtrans}
The group
$\mathrm{Gal}(K/\mathbb{Q})$ acts transitively
on the set of $k$-tuples
\[ \mathcal{S}:=\left\{(\alpha_1,\ldots,\alpha_k) \in K^k \mid 
\alpha_1^{d_1}=a_1 \wedge \cdots \wedge \alpha_k^{d_k}=a_k 
\right\} \]
 \end{lemma}
\begin{proof}
Recall that $\mysqrt{d_1}{a_1},\ldots,\mysqrt{d_k}{a_k}$ are real radicals with respective minimal polynomials $x^{d_i}-a_i$ and~$a_i$ mutually coprime.


Let $L_i := \QQ(\mysqrt{d_1}{a_1}, \ldots, \mysqrt{d_i}{a_{i}})$ for $i\in\{1,\ldots,k\}$.
By virtue of the~$a_i$ being coprime and~\cite[Lemma 4.6]{blomer98}, 
the polynomial~$f_i := x^{d_i}-a_i$ stays irreducible over~$L_{i-1}$, and thus is the minimal polynomial of $\mysqrt{d_i}{a_i}$
 over this field.
 
The proof follows by repeated use of the Isomorphism extension theorem, cf. \cite[Theorem~5.12]{StewartBook}.
Note that $L_{i}$ is a simple extension of $L_{i-1}$ with $L_{i} = L_{i-1}(\alpha_i)$ where $\alpha_i$ is a solution of $f_i$. Denote by $\psi_{i-1}$ an embedding of $L_{i-1}$ into $K$. If $\alpha_i'$ is a root of $\psi_{i-1}(f_i)$ then there is a unique extension of $\psi_{i-1}$ to a homomorphism $\psi_i : L_i \to K$ such that $\psi_i(\alpha_i) = \alpha_i'$.
Applying the above inductively, we obtain a homomorphism $\psi_{k} : \QQ(\mysqrt{d_1}{a_1}, \ldots, \mysqrt{d_k}{a_{k}}) \to K$, which by the Isomorphism extension theorem can again be extended to an automorphism of $K$ acting jointly transitively on $\mathcal{S}$.
\end{proof}


We now show that for an instance of  RIT, that is, an algebraic circuit with underlying polynomial $f$, and radical input with pairwise coprime radicands, the finite field computation is sound.
Given a rational prime~$p$ and a polynomial $g(x)\in \ZZ[x]$ we denote by $\bar{g}\in \FF_p[x]$  the reduction of $g$ modulo $p$.

\begin{lemma}\label{th:RIT-soundess}
Let $p$ be a prime that completely splits in $\OO_K$, and
let $\overline{\alpha}_1,\ldots,\overline{\alpha}_k \in \mathbb{F}_p$
be roots of the polynomials $x^{d_1}-a_1,\ldots,x^{d_k}-a_k$, respectively.
Then for all 
 $f(x_1,\ldots,x_k)\in\ZZ[x_1,\ldots,x_k]$, we have
\begin{enumerate}
	\item if $f(\mysqrt{d_1}{a_1},\ldots,\mysqrt{d_k}{a_k}) = 0$ then $\overline{f}(\overline{\alpha}_1,\ldots,\overline{\alpha}_k) = 0$, and
	\item if $\overline{f}(\overline{\alpha}_1,\ldots,\overline{\alpha}_k) = 0$ then $p\mid N_{K/\QQ}(f(\mysqrt{d_1}{a_1},\ldots,\mysqrt{d_k}{a_k}))$.
\end{enumerate}
\end{lemma}
\begin{proof}
Recall that the radicals $\mysqrt{d_1}{a_1},\ldots, \mysqrt{d_k}{a_k}$ are such that their respective minimal polynomials are $x^{d_1}-a_1,\ldots,x^{d_k}-a_k$.

Consider a prime-ideal factor $\mathfrak{p}$ of $p\OO_K$.
The quotient homomorphism $\OO_K \rightarrow \mathbb{F}_p$
with kernel $\mathfrak{p}$ maps the $d_i$ distinct roots of 
each polynomial $x^{d_i}-a_i$ in $\OO_K$ bijectively onto the 
roots of the same polynomial in $\mathbb{F}_p$.
Then, by joint transitivity of the Galois group $\mathrm{Gal}(K/\mathbb{Q})$, established in~\Cref{lem:jtrans}, there is a homomorphism
$\varphi:\OO_K\rightarrow \mathbb{F}_p$ such that 
$\varphi(\mysqrt{d_i}{a_i})=\overline{\alpha}_i$ for all $i \in \{1,\ldots,k\}$.

Item 1 follows from the fact that if $f(\mysqrt{d_1}{a_1},\ldots,\mysqrt{d_k}{a_k}) =0$
then $\overline{f}(\overline{\alpha}_1,\ldots,\overline{\alpha}_k)=\varphi(f(\sqrt[d]{a_1},\ldots,\sqrt[d]{a_k}))=0$.
For Item 2, note that
the kernel of $\varphi$ is a prime ideal of $\OO_K$ lying above $p$.
Thus if $\overline{f}(\overline{\alpha}_1,\ldots,\overline{\alpha}_k)=0$
then $p\mid N_{L/\QQ}(f(\alpha_1,\ldots,\alpha_k))$.
\end{proof}

We have just shown that given an instance of  RIT, the computation can be taken to a finite field $\FF_p$ for some rational prime $p$. In the following section, we discuss how to choose an appropriate prime~$p$ that satisfies the two conditions given in \Cref{th:RIT-soundess}.

\Cref{lem:jtrans} plays an important role in the construction of our algorithm, and furthermore is one of the properties of the input to  RIT  that makes our technique difficult to generalise to more general identity testing problems. In particular, joint transitivity ensures that in \Cref{th:RIT-soundess}(1), no matter which representative of the $\mysqrt{d_i}{a_i}$ we choose in $\FF_p$, that is, no matter which solution of the equation $x^{d_i}-a_i$ we guess in $\FF_p$, the computation remains sound. If we were, for instance, to generalise our identity testing problem to allow radical and cyclotomic inputs, joint transitivity may not hold anymore.

\begin{example}
\label{ex:generalisation}
Consider the polynomial $f(x_1,x_2) = x_2^2 - x_1 x_2 + 1$ with input $\sqrt{2}$ and a primitive $8$-th root of unity $\zeta_8$. The polynomial $f$ vanishes at $(\sqrt{2},\zeta_8)$. The number field of the computation is $\QQ(\sqrt{2},\zeta_8)$, and we can choose the completely split prime $17$ for our finite field computation. The minimal polynomials of our input split as $x^2 - 2 = (x-6)(x+6)$ and $x^4 + 1 = (x+2)(x-2)(x+8)(x-8)$ in $\FF_{17}$. However, since the Galois group of the field $\QQ(\sqrt{2},\zeta_8)$ does not act jointly-transitively on the input, we cannot choose the representatives of our two input numbers in~$\FF_{17}$ arbitrarily. In particular, by choosing $6$ for $\sqrt{2}$ and $2$ for $\zeta_8$, and evaluating~$f$ in~$\FF_{17}$, the result would be $10$, a clear false negative. This is due to the fact that the minimal polynomial $\Phi_8(x) = x^4 + 1$ of $\zeta_8$ reduces in $\QQ(\sqrt{2})$, hence, as soon as we choose $6$ as a representative for $\sqrt{2}$, we cannot choose the representative for $\zeta_8$ freely. In fact, if we replace $x_1$ by $\sqrt{2}$ and $x_2$ by $x$ in the polynomial~$f$, we obtain $(x^2 - \sqrt{2}x + 1)$,
 which is a factor of $\Phi_8(x)$ in $\QQ(\sqrt{2})$. Indeed, $\Phi_8(x)$ factors as $(x^2 - \sqrt{2}x + 1)(x^2 + \sqrt{2}x + 1)$ in $\QQ(\sqrt{2})$, hence, as soon as we choose $6$ as a representative for $\sqrt{2}$, we can only choose the roots of
 $(x^2 - 6x + 1)$
 as representatives for $\zeta_8$, whereas $2$
 is a root of the polynomial $(x^2 + 6x + 1)$ in $\FF_{17}$.
\end{example}

We note here that the underlying approach (working modulo prime ideals) is the same as in \cite{balaji2021cyclotomic} and preceding works on versions of ACIT.
However, as shown in \Cref{ex:generalisation}, the transitivity condition given in 
\Cref{lem:jtrans} is crucial in our approach, whereas it has no equivalent in \cite{balaji2021cyclotomic}.
Furthermore, we have to allow for the fact that the ring of integers of the number field is no longer monogenic in the present case (e.g., \Cref{th:RIT-soundess} is analogous to Theorem~8 in \cite{balaji2021cyclotomic}, but the proof requires working in a certain order in the ring of integers). 
In relation to this, \Cref{ex:generalisation} explains some of the difficulties arising from attempting a common generalisation of \cite{balaji2021cyclotomic} and the present paper.


\subsection{Choice of the prime $p$}

Let us call the primes $p$ that completely split in $\OO_K$
\emph{good} primes, and the primes that do not divide the norm of the algebraic integer computed by the circuit \emph{eligible} primes. We will first discuss how to choose eligible primes.
All missing proofs of Lemmas we establish in this subsection can be found in \Cref{appendix:conp}.


Given a rational prime $p$, we would like to ensure that $p \nmid N(\alpha)$, where $\alpha$ is the algebraic integer computed by a circuit.
The norm $N(\alpha)$ has at most $\log |N(\alpha)|$ prime divisors. Intuitively, this means that if we have a set of $\log |N(\alpha)| + 1$ primes, at least one of them will be eligible. Now let us see how we can bound this value.

Recall that the norm of an algebraic number $\alpha$ in a Galois field can be computed as the product of its Galois conjugates. Thus, in order to bound the magnitude of the norm of our computed number, we need a bound on the magnitude of the conjugates of $\alpha$, as well as the size of the Galois group, that is, the number of conjugates of $\alpha$. The latter can be obtained using a bound on the degree of the number field $K$ itself.
Recall that  the splitting field~$K$ of $\prod_{i=1}^k (x^{d_i}-a_i)$ is $\QQ(\mysqrt{d_1}{a_1}, \ldots, \mysqrt{d_k}{a_k},\zeta_d)$. Since the $d_i$ are of magnitude at most 
$2^s$, it follows that $d$ is of magnitude at most $2^{s^2}$,
and the degree of $K$ will be at most $2^{2s^2}$.
Using this bound, we show:

\begin{restatable}[Bound on the norm]{lemma}{boundnorm}\label{lem:norm}
Denote by $\alpha\in \OO_{K}$ the algebraic integer computed by~$C$ evaluated on the~$\mysqrt{d_i}{a_i}$. We have
\[ |N(\alpha)| \leq 2^{2^{s^3}} \]
for $s\geq 4$.

\end{restatable}

In the context of our algorithm, this means that if we find 
$2^{s^3} + 1$ good primes, at least one of them will also be eligible, and hence our finite field computation will be sound. Let us now focus on how we can ensure to be able to find enough good primes of polynomial bit-length in the size of the input to complete our reasoning.

To this aim, we use a quantitative version of the Chebotarev density theorem. Intuitively speaking, given a Galois extension $L$ of $\QQ$, the theorem gives a bound on the number of primes splitting in a certain pattern in $\OO_L$. The different classes of splitting patterns correspond to conjugacy classes of the Galois group $\Gal(L/\QQ)$ of $L$. As it turns out, the primes splitting completely in $L$ correspond to the conjugacy class $\{id\}$  
containing solely the identity element~$id$ of $\Gal(L/\QQ)$. The asymptotic version of the theorem then asserts that the set of completely split primes has density $\frac{1}{|\Gal(L/\QQ)|}$.  Denoting by $\pi(x)$ the number of all rational primes less or equal to $x$, and by $\pi_1(x)$ the number of completely split primes, the quantitative version of the theorem is as follows~\cite{lagariasO, serre1981quelques}:

\begin{proposition}[Bound on $\pi_{1}(x)$]\label{prop:number_of_primes}
Assuming GRH,
\[ \pi_{1}(x) \geq \frac{1}{|\Gal(K/\QQ)|} \left[ \pi(x) - \log \Delta_K - c x^{1/2} \log (\Delta_K x^{|\Gal(K/\QQ)|}) \right] \]
where $c$ is an effective constant.
\end{proposition}

To apply the above proposition we will need a bound on the discriminant~$\Delta_K$ of the number field $K$. Recall the definition of the discriminant of a number field; given a $\ZZ$-basis $\{\alpha_1,\ldots,\alpha_n\}$ of the ring of integers $\OO_L$ of a number field $L$, the discriminant $\Delta_L$ is the determinant of the matrix $\mathrm{tr}(\alpha_i\alpha_j)$ for all $i,j=1,\ldots,n$.
However, in the case of a radical field extension $L$, 
 computing a basis for its ring of integers~$\OO_L$ is a non-trivial task. In order to avoid this computation, we try to bound the discriminant using the discriminant of an order of our ring of integers $\OO_K$. Recall that an order~$\OO$ in a number field $L$
is a free $\ZZ$-submodule of $\OO_L$ of rank $[L:\QQ]$. 
Looking again at the number field $L=\QQ(\sqrt{5})$ with ring of integers $\OO_L=\ZZ[\frac{\sqrt{5}+1}{2}]$, note that $\ZZ[\sqrt{5}] \subset \OO_L$; in particular, $\ZZ[\sqrt{5}]$ is an order of index 2 in $\OO_L$. The following (see, e.g., \cite[Proposition 4.4.4]{cohen2013course}) holds:

 \begin{proposition}
 \label{prop:order_disc}
 Suppose $\OO$ is an order in $\OO_L$. Then
 \[ \Disc(\OO) = \Disc(\OO_L) \cdot [\OO_L : \OO]^2 \]
 \end{proposition}
 
We construct an order $\OO$ of $\OO_K$, the discriminant of which we can bound using a standard result in algebraic number theory. The Primitive element theorem states that any number field $L$ can be generated by adjoining a single element $\theta$, called the primitive element, to $\QQ$, i.e., $L = \QQ(\theta)$. Then the subring $\ZZ[\theta]$ of $\OO_L$ is an order of $\OO_L$. Furthermore, the discriminant $\Disc(\ZZ[\theta])$ is equal to the discriminant of the minimal polynomial of $\theta$.
It is well-known that the discriminant of a polynomial can be computed as the product of the differences of the roots.

We thus start with preliminary result regarding primitive elements of the number field~$K$.

\begin{restatable}[Bound on the primitive element]{lemma}{boundprimitiveelement}\label{lem:primitive_el}
The field $K$ has a primitive element $\theta$, computed as the linear combination
\[ \theta = c_0 \zeta_d + \sum_{i=1}^k c_i \mysqrt{d_i}{a_i}\]
with $c_i \leq 2^{4s^2} \in \ZZ$ and $\deg \theta \leq 2^{2s^2}$.
\end{restatable}

Henceforth, we fix a primitive element~$\theta$ for our number field~$K$, computed as in Lemma~\ref{lem:primitive_el}. Now \Cref{prop:order_disc}
suggests that $\Delta_K \leq \Delta_{f_\theta}$, where $\Delta_{f_\theta} = \Disc(\ZZ[\theta])$, which we bound as follows:


\begin{restatable}[Bound on the discriminant]{lemma}{bounddiscriminant}\label{lem:discriminant}
 We have
 \[ |\Disc(\ZZ[\theta])| \leq 2^{2^{5s^2}}\]
 for $s\geq 4$.
\end{restatable}

Recall that we would like to choose enough good primes $p$ so that at least one of them will be eligible, i.e., at least one of them will not divide the norm of the computed algebraic integer. In particular, we would like to find $x$ such that
 $\pi_{1}(x) \geq 2^{s^3} + 1$.
Using the bound above, we claim that this is the case for 
$x \geq 2^{4s^3}$:


\begin{restatable}{lemma}{boundnprimes}\label{prop:bound_p_conp}
Assuming GRH,
\[ \pi_{1}(2^{4s^3}) \geq  2^{s^3}  + 1.\]
\end{restatable}

\subsection{Proof of \Cref{th:rit-conp}}


\begin{proof}[Proof of Theorem~\ref{th:rit-conp}]
	\Cref{fig:conp_RIT} presents a nondeterministic polynomial time algorithm for the complement of  RIT  as follows.
	
	Given input radicals $\mysqrt{d_1}{a_1},\ldots,\mysqrt{d_k}{a_k}$,
	denote by $K$ the splitting field  of $\prod_{i=1}^k (x^{d_i}-a_i)$.
	 Further denote by $\theta$ a primitive element of~$K$, computed as in Lemma~\ref{lem:primitive_el}. 
	
	Let us first argue that the algorithm runs in polynomial time. 
In Step 1, after guessing candidates for $p$ such that $p \equiv 1 \pmod{d}$ and $\overline{\alpha}_1,\ldots,\overline{\alpha}_k$,  verifying whether $\overline{\alpha}_i^{d_i} \equiv a_i \pmod{p}$ can be done in polynomial time by the repeated-squaring method. It is clear that Step 2 can be done in polynomial time.
	
	Now let us show that the RIT problem is in \coNP. Suppose $f(\mysqrt{d_1}{a_1},\ldots,\mysqrt{d_k}{a_k}) \neq 0$. Under GRH, the lower bound in \Cref{prop:bound_p_conp} shows that
	 $\pi_{1}(2^{4s^3}) \geq 2^{s^3} + 1$.
	 It follows that there exists a prime
	 $p\leq 2^{4s^3}$  such that 
	\begin{itemize}
		\item $p \nmid N(f(\mysqrt{d_1}{a_1},\ldots,\mysqrt{d_k}{a_k}))$, and
		\item $p$ splits completely in $K$.
	\end{itemize}
	
	The polynomial certificate of non-zeroness of $f(\mysqrt{d_1}{a_1},\ldots,\mysqrt{d_k}{a_k})$ then comprises, the prime $p$ above, as well as the list of integers $\overline{\alpha}_1,\ldots,\overline{\alpha}_k\in\FF_p$ such that $\overline{\alpha}_i^{d_i}\equiv a_i \pmod{p}$. Following \Cref{th:RIT-soundess}, we then have that $\overline{f}(\overline{\alpha}_1,\ldots,\overline{\alpha}_k) \neq 0$.
	   
	   On the other hand, as we have noted above, for any prime~$p$ and the representation~$\overline{\alpha}_1,\ldots,\overline{\alpha}_k$ of radicals $\mysqrt{d_1}{a_1},\ldots,\mysqrt{d_k}{a_k}$ in $\FF_p$, if 
	   $f(\mysqrt{d_1}{a_1},\ldots,\mysqrt{d_k}{a_k}) = 0$, then $\overline{f}(\overline{\alpha}_1,\ldots,\overline{\alpha}_k) = 0$, as shown in \Cref{th:RIT-soundess}, which concludes the proof.
\end{proof}

